\section{Contesto}
\label{sec:background}

\subsection{FedAvg}
FedAvg aggiorna il modello globale come media dei
pesi locali, pesata per il numero di campioni di ciascun client:
\begin{equation}
  w \leftarrow \sum_{i=1}^{K} \frac{n_i}{n} w_i, \qquad n = \sum_{i=1}^{K} n_i
  \label{eq:fedavg}
\end{equation}
dove $w_i$ e $n_i$ sono rispettivamente i pesi e il numero di campioni di
addestramento del client $i$. Sotto ipotesi di dati non-i.i.d., la velocità di
convergenza di FedAvg degrada con il grado di eterogeneità tra le
distribuzioni locali.

\subsection{DiLoCo e comunicazione sparsa}
DiLoCo~\cite{douillard2023diloco} riduce la frequenza di comunicazione
alternando $H$ passi di ottimizzazione locale (\emph{inner steps}) a un
singolo passo di sincronizzazione globale, eseguito da un \emph{outer
optimizer} con momentum mantenuto centralmente. Il sistema qui presentato
adotta la stessa idea di comunicazione sparsa a $H$ passi, ma non può
adottare l'outer optimizer di DiLoCo così com'è: quest'ultimo richiede uno
stato di momentum globale coerente per tutti i client, il che presuppone
un coordinatore centrale. Il sistema sostituisce
quindi la sincronizzazione globale periodica di DiLoCo con un'aggregazione
locale asincrona basata su un push one-hop a fanout fisso tra vicini,
descritta nella Sezione~\ref{sec:design} e motivata per contrasto con le
alternative del corso nella Sezione~\ref{sec:background}E.

\subsection{Dataset: LEAF e FEMNIST}
FEMNIST è un dataset di caratteri manoscritti (cifre e
lettere, 62 classi, immagini $28\times28$ in scala di grigi) partizionato per
scrittore: ogni scrittore produce una distribuzione di stili ed etichette
propria, rendendo il dataset naturalmente non-i.i.d.\ tra client quando ogni
partizione corrisponde a scrittori diversi. Nel sistema qui presentato ogni
client riceve dati di \emph{più} scrittori LEAF, per simulare
uno scenario realistico in cui un partecipante rappresenta un piccolo gruppo
di utenti piuttosto che un singolo individuo.

\subsection{Strategie alternative non implementate}
La letteratura propone diverse estensioni di FedAvg per mitigare
l'eterogeneità dei dati: FedProx~\cite{li2020fedprox}, SCAFFOLD~\cite{karimireddy2020scaffold},
FedBN~\cite{li2021fedbn}. Il sistema qui presentato implementa esclusivamente FedAvg,
scelto come strategia di riferimento più semplice da rendere verificabile in
un contesto decentralizzato puro. Non tutte le
alternative sono equivalenti dal punto di vista architetturale: FedProx
modifica solo la funzione di costo usata in Fase B (Sezione~\ref{sec:details}B)
e potrebbe essere adottata senza alcuna modifica al buffer di aggregazione o
al protocollo gRPC; SCAFFOLD richiederebbe di trasmettere anche variabili di
controllo oltre ai pesi, aumentando la dimensione dei messaggi ma senza
introdurre un aggregatore centrale.

\subsection{Scelta del protocollo di comunicazione}
Il corso di riferimento presenta tre famiglie di disseminazione a livello
applicativo su overlay non strutturato: \emph{flooding},
\emph{random walk} e \emph{gossip epidemico}, a sua volta nelle
varianti \emph{anti-entropy} e \emph{rumor mongering}. Il sistema qui
presentato usa invece un
\textbf{push one-hop a fanout fisso $k$ con grpc} (Sezione~\ref{sec:details}B, indicato
anche come \emph{k-push}); i vicini sono scelti casualmente a ogni round,
un modello neutro di una qualunque politica di selezione reale (prossimità
di rete, similarità dei dati locali se disponibile, o un altro stato osservabile del nodo). I confronti
che seguono si riferiscono alle scale testate in questo lavoro, $N\le8$ — il
tetto imposto da AWS Learner Lab (Sezione~\ref{sec:details}F), non un limite
intrinseco del protocollo.

\textbf{Perché non flooding.} Il motivo principale è il traffico: il numero
di trasmissioni cresce con il grado medio dei vicini e, soprattutto, con il
\emph{time-to-live} — ogni hop aggiuntivo moltiplica i messaggi per il
fattore di ramificazione della rete.Il
caso limite $k=N-1$ del nostro push one-hop è di fatto equivalente a un singolo round di flooding con
\emph{time-to-live}$=1$: ogni worker raggiunge direttamente tutti i peer
noti, senza bisogno di alcun inoltro successivo. La differenza pratica tra le
due soluzioni emerge solo a \emph{time-to-live}$>1$, cioè quando servirebbe
più di un hop per coprire l'intera rete. Su reti più grandi e sparse il traffico
crescerebbe in modo moltiplicativo a ogni hop: resta quindi la soluzione più
costosa delle tre a qualunque scala.

\textbf{Il push one-hop equivale a $N$ k-random-walk simultanei.} Il singolo hop di trasporto coincide di
fatto con il primo passo di un \emph{k-random-walk} con
\emph{time-to-live}$=1$. La differenza non è tanto il \emph{time-to-live} quanto la molteplicità delle origini: un
k-random-walk parte da \emph{un solo} nodo con una query da propagare; il
nostro sistema fa ripartire lo stesso meccanismo da \emph{ogni} worker, ogni
round, in parallelo.

\textbf{Gossip epidemico.} È un'alternativa legittima: il relay multi-hop propagherebbe gli
aggiornamenti più rapidamente e più fedelmente all'originale, un vantaggio
reale su reti molto più grandi di quelle testate, dove un $k$ fisso non
basterebbe più a coprire l'intera rete in pochi round. Il costo è il
traffico (requisito del progetto)— ogni hop di relay moltiplica i messaggi per il fattore di
ramificazione, mentre il push one-hop mantiene un costo per nodo
indipendente da $N$ — quindi un compromesso esplicito tra traffico e
qualità/freschezza dell'informazione, da valutare caso per caso. La nostra soluzione, essendo applicata ad un contesto
di machine learning federato con fed avg, comunque propoaga ragionevolmente nel tempo l'informazione a tutta la rete,
anche se hop è 1, in quanto i messaggi con i nuovi pesi hanno "memoria" anche dei vecchi aggiornamenti.
Infatti il guadagno del relay è inoltre già in parte attenuato dal
dominio applicativo: grazie a FedAvg, lo stato che un worker invia è già il
risultato di aggregazioni con i round precedenti, quindi la propagazione
multi-hop avviene comunque, implicitamente, nel tempo anziché nello spazio
del singolo round. Il relay porterebbe anche un
problema standard di deduplicazione — lo stesso messaggio potrebbe
raggiungere lo stesso worker da percorsi indipendenti nello stesso round —
comunque risolvibile. Un osservazione: quando il payload scambiato è l'intero stato di un modello e non un piccolo
messaggio, il costo di un relay multi-hop cresce con il fattore di
ramificazione a ogni hop, mentre l'aggregazione locale ottiene una
propagazione senza traffico aggiuntivo dedicato al solo inoltro. In definitiva, abbiamo preferito restare entro i
vincoli di traffico richiesti esplicitamente dalla traccia di progetto;
resta un'osservazione interessante da approfondire in futuro come si
comporterebbe il gossip epidemico in un contesto con molti più nodi, e come
si confronterebbe con la soluzione qui proposta.


\textbf{Verifica empirica.} Il grafo one-hop raggiunge comunque la
connettività completa alle scale qui testate: con $k=N-1$ ogni worker
contatta \emph{tutti} i propri peer a ogni round, il limite superiore
raggiungibile in un singolo hop. I risultati (Sezione~\ref{sec:results}A,
Sezione~\ref{sec:discussion}A) mostrano che questo da solo chiude la
maggior parte del divario osservato a fanout basso (+7--8 punti di
accuratezza global-test tra $k=1$ e $k=N-1$ a $N=8$, spread quasi dimezzato);
ma anche a $k=N-1$ lo spread non scende mai sotto il 14\%
(Sezione~\ref{sec:discussion}D). Poiché $k=N-1$ è già la copertura massima
raggiungibile in un hop, un relay multi-hop non avrebbe potuto ridurre
ulteriormente questo residuo: la causa più probabile non è la portata del
protocollo di comunicazione, ma l'assenza di un meccanismo esplicito di
correzione del client drift (FedProx, SCAFFOLD, Sezione~\ref{sec:background}D)
— una limitazione indipendente dalla scelta tra one-hop e multi-hop, discussa
in Sezione~\ref{sec:discussion}D.
